\section{The program and the sufficiency of its optimality conditions}
\label{sec:problem}

Write~\eqref{eq:qp} as
\begin{equation}
  \min_{x \in \R^n} \; f(x) := \tfrac12 x\T G x - a\T x
  \qquad \text{subject to} \qquad
  \begin{cases}
    c_i\T x = b_i, & i \le m_{\mathrm{eq}},\\
    c_i\T x \ge b_i, & i > m_{\mathrm{eq}},
  \end{cases}
  \label{eq:qp2}
\end{equation}
with $G \succ 0$, and let
$\Omega = \{x : c_i\T x = b_i \ (i \le m_{\mathrm{eq}}),\ c_i\T x \ge b_i \
(i > m_{\mathrm{eq}})\}$ be the feasible set, a closed convex polyhedron. Since
$G \succ 0$ the objective is strictly convex and coercive, so~\eqref{eq:qp2} has
exactly one solution whenever $\Omega \neq \emptyset$; see~\cite[Ch.~16]{nocedal2006}.

The Lagrangian is $L(x,\lambda) = \tfrac12 x\T G x - a\T x - \lambda\T(C\T x - b)$
and the Karush--Kuhn--Tucker conditions are
\begin{align}
  G x - a &= C \lambda, \label{eq:stat}\\
  c_i\T x &= b_i \quad (i \le m_{\mathrm{eq}}), \qquad
  c_i\T x \ge b_i \quad (i > m_{\mathrm{eq}}), \label{eq:pfeas}\\
  \lambda_i &\ge 0 \quad (i > m_{\mathrm{eq}}), \label{eq:dfeas}\\
  \lambda_i\,(c_i\T x - b_i) &= 0 \quad (i > m_{\mathrm{eq}}). \label{eq:comp}
\end{align}
Multipliers of equality constraints are unrestricted in sign. What the argument of
Section~\ref{sec:certificate} needs is the converse of these necessary conditions,
which we state and prove because every certificate below rests on it.

\begin{proposition}[Sufficiency]\label{prop:sufficient}
If $(x,\lambda)$ satisfies \eqref{eq:stat}--\eqref{eq:comp} then $x$ is the
unique solution of~\eqref{eq:qp2}.
\end{proposition}

\begin{proof}
Let $\tilde x \in \Omega$. By strict convexity,
$f(\tilde x) \ge f(x) + (G x - a)\T(\tilde x - x)$, with equality only if
$\tilde x = x$. Substituting~\eqref{eq:stat},
\[
  (G x - a)\T(\tilde x - x)
  = \lambda\T\bigl(C\T \tilde x - C\T x\bigr)
  = \sum_i \lambda_i\bigl[(c_i\T\tilde x - b_i) - (c_i\T x - b_i)\bigr]
  = \sum_i \lambda_i (c_i\T \tilde x - b_i),
\]
the last step by~\eqref{eq:comp} for the inequalities and by $c_i\T x = b_i$ for
the equalities. Each surviving term is non-negative: for $i \le m_{\mathrm{eq}}$
the bracket vanishes, and for $i > m_{\mathrm{eq}}$ both factors are non-negative
by~\eqref{eq:pfeas} and~\eqref{eq:dfeas}. Hence $f(\tilde x) \ge f(x)$ with
equality only at $\tilde x = x$.
\end{proof}

Proposition~\ref{prop:sufficient} is what makes an unguaranteed method safe rather
than usually right. A point equipped with multipliers passing the four conditions
is \emph{the} answer, so a method that produces candidates may be as unprincipled
as one likes provided every candidate is tested.

\subsection{The working-set subproblem}

Fix a working set $\Active$ of size $k$, and write $A = C_\Active$ for its
normals. The \emph{working-set subproblem} holds every constraint in $\Active$ as
an equality, whatever its type in~\eqref{eq:qp2}, and ignores the rest:
\begin{equation}
  \min_{x} \; \tfrac12 x\T G x - a\T x
  \qquad \text{subject to} \qquad A\T x = b_\Active .
  \label{eq:sub}
\end{equation}

\begin{proposition}[Solution of the subproblem]\label{prop:sub}
If $A$ has full column rank $k$, problem~\eqref{eq:sub} has the unique solution
and multipliers
\begin{equation}
  u = \bigl(A\T G^{-1} A\bigr)^{-1}\bigl(b_\Active - A\T x_u\bigr),
  \qquad
  x = x_u + G^{-1} A\, u,
  \qquad x_u := G^{-1} a .
  \label{eq:subsol}
\end{equation}
\end{proposition}

\begin{proof}
Stationarity for~\eqref{eq:sub} reads $Gx - a = Au$, so
$x = x_u + G^{-1}Au$. Substituting into $A\T x = b_\Active$ gives
$A\T x_u + (A\T G^{-1}A)u = b_\Active$. The matrix $A\T G^{-1} A$ is positive
definite because $G^{-1} \succ 0$ and $A$ has full column rank, hence invertible,
which yields~\eqref{eq:subsol}; the point so constructed is the unique minimiser
by strict convexity of $f$ on the affine feasible set of~\eqref{eq:sub}.
\end{proof}

We call $A\T G^{-1} A$ the \emph{dual Hessian} of the working set. Its
invertibility is equivalent to $A$ having full column rank, and that single fact
is where the obstruction of Section~\ref{sec:obstruction} comes from. Note what
the pair $(x,u)$ of~\eqref{eq:subsol}, extended by zeros off $\Active$, does and
does not deliver: stationarity~\eqref{eq:stat} and complementarity~\eqref{eq:comp}
hold by construction, since the working-set constraints have zero slack and the
rest have zero multiplier, while the sign condition~\eqref{eq:dfeas} on
$\Active$'s inequalities and primal feasibility~\eqref{eq:pfeas} off $\Active$
need not hold at all. Every method discussed below is a rule for choosing
$\Active$ so that the two missing conditions come out right.
